Domain-specific persona prompts are a natural choice for AI creative generation tasks:
instruct the model to adopt the identity of a practitioner, and domain conventions should
follow. Yet in a controlled puzzle authoring study, a bare professional label
(\emph{``You are an artist''}) outperformed an equally sparse domain-specific label
(\emph{``You are a puzzle designer''}) by 15.6 points on a seven-dimension weighted
rubric (/45), scoring 40.3 against 24.7. The cognitive mechanism behind this reversal
is unknown: the artist persona encodes no puzzle design knowledge, yet produces
categorically superior output. Prior work treats persona prompts as holistic units;
no study has decomposed a persona into constituent cognitive traits and identified
which subset is causally responsible for a quality outcome.

We address this gap through a three-phase decomposition study. Phase~1 isolates each
of six candidate traits---experience-first thinking, surprise instinct, visual thinking,
elegance drive, rigor and verification, and systematic thinking---as a single imperative
sentence evaluated by an expert panel on the same weighted rubric. Phase~2 applies
leave-one-out ablation from the full artist persona, measuring each trait's structural
contribution. Phase~3 performs cumulative reconstruction from the null baseline, adding
traits in tier order to identify the minimum sufficient set. All conditions are evaluated
by a three-member C8 expert panel on the same seven-dimension weighted rubric ($\geq 33$
passing threshold).

Experience-first thinking (T1) is architecturally irreplaceable: removing it from the
full artist persona causes a catastrophic $-29.4$~point drop (from 41.7 to 12.3),
producing a list of domain terms with no solver mechanism rather than a degraded puzzle.
Elegance drive (T4) and rigor with verification (T5) are load-bearing: removing either
drops the score below threshold ($-11.7$ and $-12.7$, respectively). In Phase~1, visual
thinking (T3) and surprise instinct (T2) pass in isolation (35.7 and 34.7,
respectively), while elegance drive (T4, 26.0) and systematic thinking (T6, 22.3) fall
at or below the domain-specific baseline---T4 because stripping without an
experience-first target removes the solver's journey along with the clutter. Phase~3
reconstruction confirms that T1\,+\,T4\,+\,T5 is the minimum sufficient set: their
cumulative combination (R3) is the first condition to cross the passing threshold, at
35.7/45. Visual thinking, surprise instinct, and systematic thinking are quality
amplifiers whose removal the structure survives.

This paper makes three contributions. First, it provides the first trait-level
decomposition of a creative persona prompt effect, replacing a holistic persona result
with a causal account. Second, it empirically identifies a minimum sufficient trait set
(T1\,+\,T4\,+\,T5)---three sentences encoding experience-first framing, elegance drive,
and verified rigor---that crosses the quality threshold and generalizes to any creative
generation task where audience experience is the primary design constraint. Third, it
demonstrates that domain-naive artistic orientation outperforms domain-specific expertise
for creative AI generation, and explains why: domain labels activate template conventions,
while artistic orientation activates the experience-design stance that makes domain
knowledge serve the audience rather than serve itself.
